
% ════════════════════════════════════════════
% 总体分析（2.总体分析.tex）写作规则 —— 模板内嵌，AI 先读本文件再写
% ════════════════════════════════════════════
% 【结构】三段式，纯文字，不画流程图（流程图放到各问 5.X.1 的具体分析里）。
% 【第1段（引文，一句话）】总体表达题目意思，引入分析，一句话即可。
% 【第2段（分析，理清问题联系）】逐问串联：按「针对问题#，本问建立了#来解决#」的模板，把全题几个小问怎么层层递进、共享什么、差异在哪讲清楚。
% 【第3段】给出全题递进关系的总结，为后文建模铺路。
% 【禁止】不用分点；流程图一律放各问具体分析。
% ════════════════════════════════════════════
\section{总体分析}

本赛题围绕“算力约束下提升大语言模型能力”这一核心目标，要求利用真实公开数据刻画参数规模、数据规模、数据质量与领域配比同交叉熵损失之间的定量关系，并在算力约束下求解最优资源配置，四道小题层层递进、前后衔接。

针对问题一，本问建立了基于熵权法的质量综合评价模型、基于成对标准化得分差与 Kendall 协同系数的冲突消解模型、以及基于岭正则的领域配比线性混合回归模型，来解决“如何量化数据质量与配比、并建立配比—损失定量关系”的问题，其输出的质量评分 $Q$、质量基线 $Q_0$ 与配比—损失系数矩阵为后续问题提供质量与配比输入。针对问题二，本问建立了含质量缺口的广义标度律 $L(N,D,Q)=E+AN^{-\alpha}+BD^{-\beta}+C(1-Q)^{\gamma}$，在经典 Chinchilla 标度律基础上引入质量维并保证 $Q=1$ 时精确退化，解决“质量与规模、数据如何共同决定损失、能否相互替代”的问题，其输出的标度律参数直接供问题三的优化目标使用。针对问题三，本问建立了严格按赛题公式分解的“基础训练 $6ND$ $+$ 质量提升 $D[g(Q)-g(Q_0)]_+$ $+$ 长文本注意力 $\eta NDL_{\text{ctx}}$”成本模型与预算约束下的 SLSQP 联合优化模型，解析给出临界上下文长度 $L_{\text{crit}}=6/\eta$，并给出结构性转移的可操作数学定义，解决“算力有限时如何在 $N$、$D$、$Q$ 之间最优配置并识别制度切换”的问题。针对问题四，本问建立了分层 Loss—Benchmark 桥接模型、含年份固定效应的面板回归分解模型、logistic 饱和外推与结构外推双路线预测模型，并完成附件 C8 的逐任务聚合与附件 C4 的宏观证据检验，解决“如何分离规模扩张与非规模技术进步的贡献并预测能力边界”的问题，将前三问以损失为目标的结论桥接为以评测得分为度量的能力演进。

由此可见，四个问题构成一条完整的建模链条：问题一提供质量与配比这一“数据侧”输入，问题二将其纳入标度律得到统一损失函数，问题三在此损失函数与成本模型上完成算力约束下的最优配置，问题四则将损失层面的结论映射到真实评测能力并外推前沿。四问共享同一套“损失最小化”的优化主线，差异在于资源的维度与约束逐步扩充；同时四问之间以显式接口相互耦合——问题一输出的质量基线 $Q_0$ 与配比修正 $M(\bm p)$ 被问题三直接引用，问题二输出的标度律参数被问题三作为目标函数读取，问题四则用 Loss—Benchmark 桥接把前三问的损失结论折算为能力得分，形成“数据 $\to$ 损失 $\to$ 算力 $\to$ 能力”的闭环。
